\documentclass[11pt,a4paper]{article}

\usepackage[utf8]{inputenc}
\usepackage[T1]{fontenc}
\usepackage[spanish,es-noquoting]{babel}
\usepackage{lmodern}
\usepackage[margin=2.5cm]{geometry}
\usepackage{array}
\usepackage{tabularx}
\usepackage{booktabs}
\usepackage[table]{xcolor}
\usepackage{hyperref}

\definecolor{verdeuteq}{RGB}{27,117,5}
\newcolumntype{L}{>{\raggedright\arraybackslash}X}
\renewcommand{\arraystretch}{1.3}

\title{GA-PRAC-01: Comparación entre modelos de proceso de requisitos en escenarios reales}
\author{
  Grupo ANY\\[4pt]
  Aldas Zambrano Alan Andrés\\
  Naranjo Flores Anderson Jeampiere\\
  Yépez Guagua Luis Antonio\\[6pt]
  Ingeniería de Requerimientos (ISR-401), 4.º nivel, paralelo A\\
  Docente: Ing. Gleiston Cicerón Guerrero Ulloa, PhD\\
  Universidad Técnica Estatal de Quevedo
}
\date{}

\begin{document}
\maketitle

\section{Criterios de comparación}
\label{sec:criterios}

\section{Escenario 1. Extractora de aceite de palma}
\subsection{Matriz comparativa del escenario 1}

\begin{center}
\small
\begin{tabularx}{\textwidth}{|>{\raggedright\arraybackslash}p{2.6cm}|L|L|L|}
\hline
\rowcolor{verdeuteq}\textcolor{white}{\textbf{Criterio}} & \textcolor{white}{\textbf{Secuencial (cascada)}} & \textcolor{white}{\textbf{Iterativo e incremental}} & \textcolor{white}{\textbf{Ágil}} \\ \hline
K1. Estabilidad y cambios & Las reglas de seguridad las definió el fabricante y se mantendrán siempre igual, por eso encajan con cascada. En cambio, el proveedor de sensores se elige en el tercer mes, así que ese requisito no se sabe desde el inicio y no encaja con cascada, que necesita establecer todo al inicio. & Iterativo encaja porque, como no se sabe qué sensor se va a usar, ese requisito se puede cambiar más adelante. & Ágil podría encajar en los reportes y las pantallas, porque aún no se sabe lo que se quiere y eso se puede cambiar después. Pero lo de seguridad no cambia, así que ágil no encaja para todo el sistema. \\ \hline
K2. Partes interesadas y disponibilidad & & & \\ \hline
K3. Tipo de sistema y criticidad & & & \\ \hline
K4. Riesgo principal & & & \\ \hline
K5. Documentación y marco normativo o contractual & & & \\ \hline
\end{tabularx}
\end{center}

\subsection{Recomendación y riesgos del escenario 1}

\begin{center}
\small
\begin{tabularx}{\textwidth}{|>{\raggedright\arraybackslash}p{4.5cm}|L|}
\hline
\rowcolor{verdeuteq}\textcolor{white}{\textbf{Elemento}} & \textcolor{white}{\textbf{Respuesta}} \\ \hline
Modelo recomendado (puro o híbrido; si es híbrido, qué parte sigue cada modelo) & \\ \hline
Razón 1 (dato del escenario, criterio y consecuencia) & \\ \hline
Razón 2 & \\ \hline
Razón 3 & \\ \hline
Por qué se descartan los otros dos modelos & \\ \hline
Tensión principal del escenario (el dato que apunta a otro modelo) & \\ \hline
Riesgo 1 del modelo elegido y su mitigación & \\ \hline
Riesgo 2 del modelo elegido y su mitigación & \\ \hline
\end{tabularx}
\end{center}


\section{Escenario 2. Aplicación móvil de pedidos de comida típica}
\subsection{Matriz comparativa del escenario 2}

\begin{center}
\small
\begin{tabularx}{\textwidth}{|>{\raggedright\arraybackslash}p{2.6cm}|L|L|L|}
\hline
\rowcolor{verdeuteq}\textcolor{white}{\textbf{Criterio}} & \textcolor{white}{\textbf{Secuencial (cascada)}} & \textcolor{white}{\textbf{Iterativo e incremental}} & \textcolor{white}{\textbf{Ágil}} \\ \hline
K1. Estabilidad y cambios
& Si Rosa cambia de puntos a promociones después de aprobar el alcance, el equipo tendría que revisar los requisitos y el diseño afectados.
& Una primera versión de pedidos permitiría a Rosa probar el flujo y decidir en el siguiente ciclo si incluye promociones.
& Rosa puede reordenar las ideas sobre puntos y promociones al inicio de cada ciclo, según lo que más le sirva al negocio.
\\ \hline
K2. Partes interesadas y disponibilidad
& Rosa puede aprobar los requisitos al inicio y revisar el producto al final. Si sus observaciones llegan tarde, corregir el flujo puede exigir rehacer trabajo.
& Como Rosa está disponible casi a diario, puede probar cada versión y aclarar dudas antes del siguiente ciclo.
& Rosa puede decidir las prioridades con frecuencia y su hijo aportar ideas. Con dos desarrolladores a medio tiempo, centralizar esa decisión ayuda a evitar instrucciones contradictorias.
\\ \hline
K3. Tipo de sistema y criticidad
& Conviene definir desde el inicio cómo se protegen las direcciones y cómo se concilian las transferencias, porque un error puede exponer datos o dejar un pedido sin pagar.
& Una versión temprana permite probar cómo se vincula una transferencia con su pedido y corregir el flujo antes de añadir otras funciones.
& En el ciclo que incorpora transferencias, el equipo puede comprobar que un pedido solo aparezca como pagado después de conciliar el pago.
\\ \hline
K4. Riesgo principal
& Con 14 semanas y dos desarrolladores a medio tiempo, detectar tarde un fallo de conexión deja poco margen para corregirlo.
& Probar una versión temprana en un Android de gama baja con cortes de conexión permite ajustar el flujo de pedidos antes de ampliarlo.
& Con presupuesto fijo, el equipo puede priorizar menú, pedidos y transferencias. Los puntos y el pago con tarjeta pueden quedar para después.
\\ \hline
K5. Documentación y marco normativo o contractual
& Aunque el acuerdo de precio fijo sea informal, dejar por escrito el alcance y los criterios de aceptación permite revisar el impacto de los cambios.
& En cada versión se pueden registrar las decisiones de Rosa y los resultados de las pruebas de conciliación.
& En cada ciclo se pueden actualizar los requisitos sobre datos de clientes y guardar la evidencia de las pruebas de pago.
\\ \hline
\end{tabularx}
\end{center}

\subsection{Recomendación y riesgos del escenario 2}

\begin{center}
\small
\begin{tabularx}{\textwidth}{|>{\raggedright\arraybackslash}p{4.5cm}|L|}
\hline
\rowcolor{verdeuteq}\textcolor{white}{\textbf{Elemento}} & \textcolor{white}{\textbf{Respuesta}} \\ \hline
Modelo recomendado (puro o híbrido; si es híbrido, qué parte sigue cada modelo) & \\ \hline
Razón 1 (dato del escenario, criterio y consecuencia) & \\ \hline
Razón 2 & \\ \hline
Razón 3 & \\ \hline
Por qué se descartan los otros dos modelos & \\ \hline
Tensión principal del escenario (el dato que apunta a otro modelo) & \\ \hline
Riesgo 1 del modelo elegido y su mitigación & \\ \hline
Riesgo 2 del modelo elegido y su mitigación & \\ \hline
\end{tabularx}
\end{center}


\section{Escenario 3. Trámites municipales por contratación pública}
\subsection{Matriz comparativa del escenario 3}

\begin{center}
\footnotesize
\begin{tabularx}{\textwidth}{|>{\raggedright\arraybackslash}p{2.6cm}|L|L|L|}
\hline
\rowcolor{verdeuteq}
\textcolor{white}{\textbf{Criterio}} &
\textcolor{white}{\textbf{Secuencial (cascada)}} &
\textcolor{white}{\textbf{Iterativo e incremental}} &
\textcolor{white}{\textbf{Ágil}} \\ \hline
K1. Estabilidad y cambios &
La ERS debe aprobarse en el mes 2. Los cambios del flujo de patentes y la posible ordenanza obligarían a tramitar modificaciones formales después de esa aprobación. &
Permite revisar en ciclos el flujo de patentes y adaptar el cálculo de tasas si cambia la ordenanza, sin perder el control de la versión aprobada. &
Ayudaría a ajustar el flujo y las tasas con rapidez, pero los cambios de alcance siguen sujetos a trámites contractuales que pueden durar semanas. \\ \hline
K2. Partes interesadas y disponibilidad &
Las aprobaciones por fase se pueden coordinar con el administrador del contrato, que es el único canal oficial y se reúne semanalmente. &
Las revisiones de cada incremento pueden programarse con el administrador; para validar con ventanilla y rentas, él debe convocar a esos usuarios. &
La priorización frecuente se complica: el administrador atiende una vez por semana y los usuarios finales no pueden participar libremente en cada ciclo. \\ \hline
K3. Tipo de sistema y criticidad &
Al manejar cobros y datos ciudadanos, una especificación formal ayuda a documentar reglas y controles antes de construir el sistema web. &
Permite probar por partes los trámites, las tasas y el acceso a datos, detectando errores antes de la instalación prevista para el mes 6. &
Las entregas cortas permitirían mostrar pantallas, pero no bastaría una lista de tareas sin pruebas y registros formales sobre cobros y datos. \\ \hline
K4. Riesgo principal &
La base catastral antigua no tiene documentación. Si su integración se analiza tarde, puede comprometer la entrega del sistema en el mes 6. &
Un incremento técnico temprano puede comprobar cómo leer los datos catastrales y reducir el riesgo de descubrir incompatibilidades al final. &
Un prototipo temprano también ayudaría, aunque las decisiones sobre acceso y cambios deben pasar por el canal contractual del GAD. \\ \hline
K5. Documentación y contrato &
Se ajusta a los hitos pactados: ERS aprobada en el mes 2, sistema probado en el mes 6 y acta de entrega en el mes 7. &
Puede mantener la ERS y la trazabilidad actualizadas en cada ciclo y conservar las aprobaciones exigidas en los pliegos. &
El backlog resulta útil para organizar tareas, pero no reemplaza la ERS aprobada ni las actas y pruebas exigidas para cobrar los hitos. \\ \hline
\end{tabularx}
\end{center}

\subsection{Recomendación y riesgos del escenario 3}

\begin{center}
\small
\begin{tabularx}{\textwidth}{|>{\raggedright\arraybackslash}p{4.5cm}|L|}
\hline
\rowcolor{verdeuteq}
\textcolor{white}{\textbf{Elemento}} & \textcolor{white}{\textbf{Respuesta}} \\ \hline
Modelo recomendado (puro o híbrido; si es híbrido, qué parte sigue cada modelo) &
Se propone un enfoque híbrido: control secuencial para aprobar la ERS, gestionar cambios del contrato y cerrar los hitos; trabajo iterativo e incremental para comprobar la integración catastral y desarrollar los módulos de trámites y cobros. \\ \hline
Razón 1 (dato del escenario, criterio y consecuencia) &
K5: el contrato exige una ERS aprobada en el mes 2 y tres entregables con pago asociado. Por ello, conviene una línea base formal que permita demostrar qué se acordó en cada hito. \\ \hline
Razón 2 &
K1: la jefa de rentas ya cambió dos veces el flujo de patentes y se prevé una posible ordenanza en el mes 4. Trabajar por ciclos ayuda a revisar esas partes sin rehacer todo el sistema. \\ \hline
Razón 3 &
K4: la base catastral antigua carece de documentación. Probar la integración en un incremento temprano reduce la posibilidad de descubrir un bloqueo técnico cerca de la entrega del mes 6. \\ \hline
Por qué se descartan los otros dos modelos &
No se recomienda cascada pura: congelar todo en el mes 2 dejaría poco margen ante los cambios de patentes, tasas e integración. Tampoco se recomienda ágil puro: el cliente no participa a diario y el contrato exige documentos, aprobaciones y trámites formales para cambiar el alcance. \\ \hline
Tensión principal del escenario (el dato que apunta a otro modelo) &
El contrato empuja hacia la planificación secuencial, pero las modificaciones del flujo de patentes, la posible ordenanza y la base catastral desconocida requieren revisión iterativa. \\ \hline
Riesgo 1 del modelo elegido y su mitigación &
Riesgo: que cada incremento introduzca cambios que contradigan la ERS aprobada. Mitigación: mantener una matriz de trazabilidad y registrar con el administrador cada solicitud de cambio antes de modificar el alcance contractual. \\ \hline
Riesgo 2 del modelo elegido y su mitigación &
Riesgo: que la integración con el catastro retrase las pruebas del mes 6. Mitigación: hacer una prueba técnica de lectura de datos al comienzo, documentar los campos disponibles y reservar tiempo para corregir incompatibilidades. \\ \hline
\end{tabularx}
\end{center}


\section{Escenario 4. Plataforma web de la cooperativa de cacao}
\subsection{Matriz comparativa del escenario 4}

\begin{center}
\small
\begin{tabularx}{\textwidth}{|>{\raggedright\arraybackslash}p{2.6cm}|L|L|L|}
\hline
\rowcolor{verdeuteq}
\textcolor{white}{\textbf{Criterio}}
& \textcolor{white}{\textbf{Secuencial (cascada)}}
& \textcolor{white}{\textbf{Iterativo e incremental}}
& \textcolor{white}{\textbf{Ágil}} \\ \hline
K1. Estabilidad y cambios
& Las reglas de liquidación del estatuto pueden definirse al inicio. Un cambio posterior en la trazabilidad solicitada por un exportador exigiría revisar el alcance aprobado.
& Cada ciclo puede ajustar la trazabilidad según la temporada, manteniendo las reglas de liquidación ya aprobadas.
& Las exigencias de los exportadores pueden entrar en la lista de prioridades del siguiente ciclo. Los cambios importantes siguen sujetos a la aprobación de la directiva.
\\ \hline
K2. Partes interesadas y disponibilidad
& La definición inicial debe coordinarse con la reunión mensual de la directiva y con los días de acopio, cuando se puede consultar a los socios.
& La gerente puede revisar versiones casi a diario. El equipo puede probarlas con socios en días de acopio y llevar las decisiones importantes a la reunión mensual.
& La gerente puede resolver dudas frecuentes, pero una decisión reservada a la directiva puede quedar pendiente durante un ciclo corto.
\\ \hline
K3. Tipo de sistema y criticidad
& Como la plataforma maneja pagos, créditos y datos personales de los socios, conviene acordar temprano las reglas de cálculo y los criterios de revisión.
& Una versión de prueba permite verificar liquidaciones con ejemplos de peso, humedad y calidad antes de usar el cálculo con pagos reales.
& Las funciones pueden entregarse por partes. El incremento que calcula liquidaciones debe pasar las pruebas del estatuto antes de usarse para pagar a los socios.
\\ \hline
K4. Riesgo principal
& Si se descubre al final que se pierden entregas al recuperar la conexión, habría que corregir el registro y revisar las liquidaciones que dependen de él.
& Un ciclo temprano permite probar el registro sin conexión en el centro de acopio y comprobar que cada entrega se sincronice una sola vez.
& El equipo puede priorizar las pruebas de pérdida y duplicación de entregas por encima de los avisos y convocatorias.
\\ \hline
K5. Documentación y marco normativo o contractual
& El estatuto sirve para documentar las reglas de liquidación desde el inicio. El plan debe prever funciones terminadas para los informes semestrales del donante.
& Cada ciclo deja decisiones aprobadas, resultados de pruebas y evidencia de funciones terminadas que se pueden incorporar al informe semestral.
& En cada ciclo se pueden registrar las exigencias de trazabilidad del comprador, su aprobación y las pruebas correspondientes para conservar el historial de cambios.
\\ \hline
\end{tabularx}
\end{center}

\subsection{Recomendación y riesgos del escenario 4}

\begin{center}
\small
\begin{tabularx}{\textwidth}{|>{\raggedright\arraybackslash}p{4.5cm}|L|}
\hline
\rowcolor{verdeuteq}
\textcolor{white}{\textbf{Elemento}}
& \textcolor{white}{\textbf{Respuesta}} \\ \hline
Modelo recomendado
& Iterativo e incremental con aprobación formal de las reglas de liquidación. La gerente revisa las versiones y la directiva aprueba los cambios importantes.
\\ \hline
Razón 1
& K1. Las liquidaciones siguen reglas estables del estatuto, mientras que los exportadores cambian las exigencias de trazabilidad por temporada. Los ciclos permiten ajustar la trazabilidad sin reabrir las reglas de cálculo.
\\ \hline
Razón 2
& K2. La gerente está disponible casi a diario y la directiva se reúne una vez al mes. El equipo puede revisar avances con la gerente y preparar las decisiones para la reunión mensual.
\\ \hline
Razón 3
& K4. La conexión en el centro de acopio es inestable. Probar el registro sin conexión en un ciclo temprano permite detectar pérdidas o duplicados antes de usar esos datos en las liquidaciones.
\\ \hline
Por qué se descartan los otros dos modelos
& Un plan en cascada con alcance fijado al inicio requeriría gestionar cambios cada vez que los exportadores actualicen la trazabilidad. Un enfoque ágil con repriorización frecuente dependería de decisiones que la gerente no puede tomar sola cuando el cambio es importante.
\\ \hline
Tensión principal del escenario
& Las reglas estables del estatuto y el impacto de los errores en los pagos favorecen una definición y aprobación inicial, como en cascada. Esa base aprobada puede implementarse y verificarse por ciclos.
\\ \hline
Riesgo 1 del modelo elegido y su mitigación
& Retraso por decisiones pendientes de la directiva. Preparar antes de cada reunión mensual las propuestas y su impacto, y avanzar entre reuniones con las tareas que ya tienen aprobación.
\\ \hline
Riesgo 2 del modelo elegido y su mitigación
& Pérdida o duplicación de entregas al sincronizar. Asignar un identificador único a cada entrega y probar cortes, reconexiones y reintentos antes de usar los registros para calcular pagos.
\\ \hline
\end{tabularx}
\end{center}

\section{Actores y framework de IR en el escenario elegido}
\subsection{Actores del proceso}

\begin{center}
\small
\begin{tabularx}{\textwidth}{|L|L|L|}
\hline
\rowcolor{verdeuteq}\textcolor{white}{\textbf{Actor del escenario}} & \textcolor{white}{\textbf{Rol en el proceso de requisitos}} & \textcolor{white}{\textbf{Cuándo y cómo participa}} \\ \hline
 & & \\ \hline
 & & \\ \hline
 & & \\ \hline
\end{tabularx}
\end{center}

\subsection{Framework de IR por fase o iteración}

\begin{center}
\small
\begin{tabularx}{\textwidth}{|L|L|L|L|}
\hline
\rowcolor{verdeuteq}\textcolor{white}{\textbf{Fase o iteración}} & \textcolor{white}{\textbf{Actividades centrales (elicitación, documentación, negociación)}} & \textcolor{white}{\textbf{Actividades transversales (validación, gestión)}} & \textcolor{white}{\textbf{Artefactos producidos}} \\ \hline
 & & & \\ \hline
 & & & \\ \hline
 & & & \\ \hline
\end{tabularx}
\end{center}


\section{Declaración de uso de IA}

% \bibliographystyle{IEEEtran}
% \bibliography{referencias}

\end{document}
